% !TeX program = lualatex
\documentclass[12pt, a4paper]{article}

% Paquetes de idioma
\usepackage[spanish,es-tabla,es-nosectiondot]{babel}
\usepackage{enumitem}
\usepackage{array}
\usepackage{booktabs}
\usepackage{tabularx}
\usepackage{xltabular}
\usepackage{colortbl}
\usepackage{upquote}
\usepackage{fancyvrb}
\usepackage{graphicx}

% ------------------------------------------------------------
%  Tipografía: LuxSans (UNL), igual que InformesPFC/Informe_Final.
%  Compilar con LuaLaTeX. Los .ttf viven en ./fonts/
% ------------------------------------------------------------
\usepackage{fontspec}
\defaultfontfeatures{Ligatures=TeX}
% LuxSans es un subconjunto de Lato: los glifos que no trae se toman de Lato
\directlua{
  luaotfload.add_fallback("luxfallback", {
    "file:Lato-Regular.ttf:mode=node;script=latn;",
  })
}
\defaultfontfeatures+{RawFeature={fallback=luxfallback}}
\setmainfont{LuxSans}[
  Path             = fonts/,
  Extension        = .ttf,
  UprightFont      = *-Regular,
  BoldFont         = *-Bold,
  ItalicFont       = *-Regular,
  ItalicFeatures   = {FakeSlant=0.2},
  BoldItalicFont   = *-Bold,
  BoldItalicFeatures = {FakeSlant=0.2},
  FontFace         = {l}{n}{Font = *-Light},
  FontFace         = {k}{n}{Font = *-Black},
]
\setsansfont{LuxSans}[
  Path             = fonts/,
  Extension        = .ttf,
  UprightFont      = *-Regular,
  BoldFont         = *-Bold,
  ItalicFont       = *-Regular,
  ItalicFeatures   = {FakeSlant=0.2},
  BoldItalicFont   = *-Bold,
  BoldItalicFeatures = {FakeSlant=0.2},
]
\usepackage{setspace}
\setlength{\emergencystretch}{2em}
\usepackage{titlesec}

% Márgenes
\usepackage[a4paper,top=3cm,bottom=3cm,left=3cm,right=3cm]{geometry}

\usepackage[colorlinks=true, linkcolor=black, citecolor=blue, urlcolor=blue]{hyperref}

% Colores
\usepackage[table]{xcolor}
\definecolor{azulFICH}{RGB}{0, 99, 175}
\definecolor{grisClaro}{RGB}{230, 230, 230}
\definecolor{fondoNota}{RGB}{242, 246, 250}
% Mayúsculas con letras espaciadas, como en la portada del informe
\newcommand{\espaciadas}[2][18]{{\addfontfeatures{LetterSpace=#1}\MakeUppercase{#2}}}

% Legibilidad
\setstretch{1.1}
\setlength{\parskip}{6pt}
\setlist{topsep=2pt, itemsep=6pt, parsep=0pt, partopsep=0pt}
\setlength{\parindent}{1em}
\titlespacing{\section}{0pt}{16pt}{0pt}
\titlespacing{\subsection}{0pt}{12pt}{0pt}

% Recuadro para notas
\newcommand{\nota}[1]{%
  \par\noindent\colorbox{fondoNota}{\parbox{\dimexpr\linewidth-2\fboxsep\relax}{\setlength{\parindent}{0pt}#1}}\par}

% Captura de pantalla centrada con borde fino
\newcommand{\captura}[2]{\par\smallskip{\centering\setlength{\fboxsep}{0pt}\fbox{\includegraphics[width=#2\linewidth]{imagenes/guia/#1}}\par}\smallskip}

% Ruta copiable
\DefineVerbatimEnvironment{ruta}{Verbatim}{fontsize=\small, frame=single, framesep=3mm}

% ============================================================
\begin{document}
% ============================================================

\begin{center}
  \setstretch{1.2}
  {\small \espaciadas{ITZA Ingeniería} \par}
  \vspace{0.4cm}
  {\large \color{azulFICH} \espaciadas[12]{Instalar el MCP de ITZA\ para Civil 3D} \par}
  \vspace{0.3cm}
  {\small Guía paso a paso --- civil3D-mcp-itza 1.2.1 \par}
\end{center}
\vspace{0.4cm}

% ================================================================
\section*{Instrucciones para la IA que lea este documento}

\nota{Si sos una IA y el usuario te pasó este PDF, guialo en la instalación:
\begin{itemize}
  \item Explicá \textbf{un solo paso por vez} y esperá a que el usuario confirme antes de seguir.
  \item En cada paso decile qué tiene que ver para saber que salió bien.
  \item Si algo no coincide, decíselo al usuario. No inventes pasos ni rutas que no estén acá.
\end{itemize}}

% ================================================================
\section*{Paso 1 --- Descomprimí el zip}

Clic derecho sobre \texttt{civil3D-mcp-itza-kit-1.2.1.zip} $\rightarrow$ \textbf{Extraer todo}
$\rightarrow$ \textbf{Extraer}.


% ================================================================
\section*{Paso 2 --- Poné el plugin en Civil 3D}

\begin{enumerate}
  \item \textbf{Cerrá Civil 3D} si está abierto.
  \item Apretá \textbf{Windows + R}, pegá esto y Enter:
\begin{ruta}
%APPDATA%\Autodesk\ApplicationPlugins
\end{ruta}
  \captura{paso2-ejecutar.png}{0.5}
  \item Copiá ahí la carpeta \texttt{CivilMcpItza.bundle} \textbf{entera} (no solo lo que
  tiene adentro).
\end{enumerate}

\textbf{Cómo saber que salió bien:} dentro de \texttt{ApplicationPlugins} ves la carpeta
\texttt{CivilMcpItza.bundle}, y adentro de ella un archivo \texttt{PackageContents.xml} y una
carpeta \texttt{Contents}. Puede haber otras carpetas \texttt{.bundle} de Autodesk: no las toques.

\captura{paso2-applicationplugins.png}{0.95}

% ================================================================
\section*{Paso 3 --- Instalá la extensión en Claude Desktop}

\begin{enumerate}
  \item Abrí Claude Desktop. Abajo a la izquierda, tocá tu nombre $\rightarrow$
  \textbf{Configuración}.
  \captura{paso3-menu-configuracion.png}{0.35}
  \item Andá a \textbf{Extensiones} $\rightarrow$ \textbf{Configuración avanzada}.
  \item Bajá hasta \textbf{Desarrollador de extensiones}. Aparece un cartel rojo de advertencia:
  es normal. Tocá el botón \textbf{Instalar extensión}.
  \captura{paso3-instalar-extension.png}{0.95}
  \item Elegí el archivo \texttt{civil3d-mcp-itza-1.2.1.mcpb} del paso 1.
  \item Se abre la pantalla de \textbf{Civil 3D MCP (ITZA)} con otro cartel rojo (aparece con
  toda extensión que no es de Anthropic). Tocá \textbf{Instalar}.
  \captura{paso3-pantalla-instalar.png}{0.6}
  \item Se abre la ventana \textbf{Configurar}. El botón \textbf{Guardar} está gris hasta que
  cambies algo: bajá hasta \textbf{Local HTTP bridge port}, cambiá \texttt{3000} por
  \texttt{3001} y tocá \textbf{Guardar}. No toques los demás campos.
  \captura{paso3-configurar.png}{0.55}
  \item Volvé a \textbf{Extensiones}. En \textbf{Instalado en tu computadora} aparece
  \textbf{Civil 3D MCP (ITZA)}: tocá \textbf{Configurar}.
  \captura{paso3-lista-extensiones.png}{0.95}
  \item \textbf{Asegurate de que diga Habilitada}, con el interruptor prendido (en azul). Puede
  quedar apagada después de instalarla: si está apagada, prendelo.
  \captura{paso3-habilitada.png}{0.95}
\end{enumerate}

\textbf{Cómo saber que salió bien:} \textbf{Civil 3D MCP (ITZA)} figura en
\textbf{Instalado en tu computadora} y adentro dice \textbf{Habilitada}.

% ================================================================
\section*{Paso 4 --- Abrí Civil 3D}

\begin{enumerate}
  \item Abrí Civil 3D.
  \item \textbf{Solo la primera vez} aparece el cartel \textbf{«Seguridad - Archivo ejecutable
  no firmado»} por \texttt{Civil3DMcpPlugin.dll}. Tocá \textbf{Cargar siempre}.
\end{enumerate}

\textbf{Cómo saber que salió bien:} en la línea de comandos de Civil 3D escribí
\texttt{C3DMCPSTATUS} y Enter. Tiene que responder algo que empiece con
\texttt{Civil3D MCP listener running: True}.

De acá en adelante el plugin se carga solo cada vez que abrís Civil 3D.

% ================================================================
\section*{Paso 5 --- Probá}

Con Civil 3D abierto, en Claude Desktop abrí un chat nuevo y pedile:

\nota{¿Qué dibujo tengo abierto en Civil 3D y en qué unidades está?}

Claude te pide permiso para usar la extensión: aceptá.

\textbf{Cómo saber que salió bien:} responde con el nombre real del dibujo (por ejemplo
\texttt{Dibujo1.dwg}) y dice que está en metros.

Después probá algo que modifique el dibujo, por ejemplo:

\nota{Creá una alineación recta de 100 m que arranque en (0,0) y vaya hacia el este.}

\end{document}
